\resizebox{\linewidth}{!}{\begin{tikzpicture}[x=1cm,y=1cm,>=Latex,font=\small]
\draw (0.0,.45) circle (.12);\draw (0.0,.33)--(0.0,-.1) (-0.23,.17)--(0.23,.17) (0.0,-.1)--(-0.2,-.37) (0.0,-.1)--(0.2,-.37);
\node[above,font=\small] at (0.0,.7) {经理};
\draw[dashed,gray] (0.0,-.6)--(0.0,-7.4);
\node[draw,minimum width=2.85cm,minimum height=.9cm,align=center] at (3.3,0) {\scriptsize\textless\textless boundary\textgreater\textgreater\\排班界面};
\draw[dashed,gray] (3.3,-.6)--(3.3,-7.4);
\node[draw,minimum width=2.85cm,minimum height=.9cm,align=center] at (6.6,0) {\scriptsize\textless\textless control\textgreater\textgreater\\排班控制};
\draw[dashed,gray] (6.6,-.6)--(6.6,-7.4);
\node[draw,minimum width=2.85cm,minimum height=.9cm,align=center] at (9.899999999999999,0) {\scriptsize\textless\textless entity\textgreater\textgreater\\班次};
\draw[dashed,gray] (9.899999999999999,-.6)--(9.899999999999999,-7.4);
\node[draw,minimum width=2.85cm,minimum height=.9cm,align=center] at (13.2,0) {\scriptsize\textless\textless entity\textgreater\textgreater\\打卡记录};
\draw[dashed,gray] (13.2,-.6)--(13.2,-7.4);
\draw[->] (0.0,-1.1)--(3.3,-1.1);
\node[above,fill=white,inner sep=2pt,font=\scriptsize,text width=3.12cm,align=center] at (1.65,-1.07) {1 提交人员、日期、门店和起止};
\draw[->] (3.3,-2.1500000000000004)--(6.6,-2.1500000000000004);
\node[above,fill=white,inner sep=2pt,font=\scriptsize,text width=3.12cm,align=center] at (4.949999999999999,-2.1200000000000006) {2 保存班次()};
\draw[->] (6.6,-3.2)--(9.899999999999999,-3.2);
\node[above,fill=white,inner sep=2pt,font=\scriptsize,text width=3.12cm,align=center] at (8.25,-3.1700000000000004) {3 验证人员日期唯一与时间范围};
\draw[->] (6.6,-4.25)--(13.2,-4.25);
\node[above,fill=white,inner sep=2pt,font=\scriptsize,text width=6.42cm,align=center] at (9.899999999999999,-4.22) {4 校验原目标与新目标均未打卡};
\draw[->] (6.6,-5.300000000000001)--(9.899999999999999,-5.300000000000001);
\node[above,fill=white,inner sep=2pt,font=\scriptsize,text width=3.12cm,align=center] at (8.25,-5.2700000000000005) {5 保存班次并记录审计};
\draw[dashed,->] (6.6,-6.35)--(3.3,-6.35);
\node[above,fill=white,inner sep=2pt,font=\scriptsize,text width=3.12cm,align=center] at (4.949999999999999,-6.319999999999999) {6 返回班次和同步结果};
\end{tikzpicture}}
